\section{Antecedentes}

El desarrollo de tecnologías de secado de café ha evolucionado significativamente en las últimas décadas, pasando de métodos tradicionales dependientes de las condiciones climáticas a sistemas mecánicos y controlados que buscan mejorar la eficiencia del proceso y preservar la calidad del producto. En este contexto, la investigación reciente ha centrado su atención en el control de variables críticas como la humedad relativa, la temperatura y el flujo de aire.

\subsection{Estado del arte de tecnologías de secado de café}

\begin{table}[htbp]
\centering
\caption{Estado del arte de tecnologías de secado de café}
\begin{tabular}{p{3cm} p{3cm} p{3cm} p{2cm} p{4cm}}
\hline
\textbf{Autor y año} & \textbf{Tecnología} & \textbf{Variables} & \textbf{Variedad (proceso)} & \textbf{Resultados relevantes} \\
\hline

Phitakwinai et al. (2019) \cite{phitakwinai_thinlayer_2019} & Secado en capa delgada con control de HR & Temperatura y humedad relativa & Arábica (lavado) & El modelo de Midilli modificado describe la cinética de secado; la HR influye en la difusividad. \\

Borém et al. (2018) \cite{borem_quality_2018} & Secador convectivo con control de punto de rocío & Temperatura y HR & Arábica (natural) & Altas temperaturas deterioran la calidad; HR moderadas preservan la integridad celular. \\

Dzaky et al. (2022) \cite{dzaky_activation_2022} & Bomba de calor con doble condensador & Temperatura y humedad específica & Robusta (lavado) & La deshumidificación mejora la eficiencia energética del proceso. \\

Abreu et al. (2025) \cite{abreu_influence_2025} & Secador mecánico de alta eficiencia & Compuestos fenólicos y calidad sensorial & Arábica (natural) & Mejora en calidad fisicoquímica y puntajes SCA cercanos a 81. \\

Fauzi et al. (2024) \cite{fauzi_influence_2024} & Secado por bomba de calor (HPD) & Temperatura y humedad del aire & Robusta (lavado) & Alta sensibilidad del proceso a la humedad del aire y reducción del consumo energético. \\

Konopatzki et al. (2019) \cite{konopatzki_price_2019} & Deshumidificación de aire (UTA) & Calidad sensorial y precio & Arábica (lavado) & Incremento en calidad y aumento del 12.11\% en el valor comercial. \\

Parra-Coronado et al. (2020) \cite{parra-coronado_preliminary_2020} & Cambios cíclicos de presión (CPCD) & Tasa de secado y compuestos volátiles & Arábica (natural) & Mejora en calidad sensorial hasta 87 puntos SCA. \\

Syamsiana et al. (2025) \cite{syamsiana_enhancing_2025} & Secador rotatorio con control difuso & Temperatura y velocidad del aire & Arábica (Honey) & Control preciso del proceso y preservación de antioxidantes. \\

\hline
\end{tabular}
\end{table}